\begin{abstract}
Interactive language agents are often improved or evaluated with the same simulated partners that provide feedback, selection, or reward. In pragmatic communication, this practice can conflate robust listener-adapted language with conventions that only work for a particular partner. We study this issue in PRAG-CrossPlay, a controlled situated reference-game benchmark in which an API speaker generates candidate messages and a listener must identify a hidden target object from a text-rendered scene. Because each listener response is scored against a ground-truth object ID, the task provides an objective communicative-success signal without an LLM judge. We compare mirror self-play, which selects messages against one fixed simulated listener, with population-play, which selects messages against a small heterogeneous listener population, and then evaluate the selected messages with held-out API listeners. Mirror self-play achieves perfect same-partner success but drops under cross-play, especially on perspective-sensitive references. In a 50-scene perspective-stress replication, mirror self-play obtains 1.000 same-play success but only 0.813 cross-play success, while population-play reaches 1.000 cross-play success on the same candidate sets; with a different held-out model family, mirror drops further to 0.713 while population-play remains at 1.000. A no-coordinate ablation shows that the strongest population-play result depends on explicit listener-invariant fallbacks and that selector quality matters when those fallbacks are removed. The results support cross-play as a simple diagnostic for distinguishing generally pragmatic messages from partner-specific communicative shortcuts.
\end{abstract}

\section{Introduction}

Pragmatic communication is not only a property of an utterance in isolation. A speaker must choose what to say with respect to a listener's perceptual context, likely alternatives, prior conventions, and possible misunderstandings. The same phrase can be sufficient for one listener and ambiguous for another: ``the small green sphere'' may work if there is only one plausible green sphere in the listener's view, but fail when a second one is present; ``the cube on the left'' may work for a partner who shares the speaker's coordinate frame, but become brittle when the listener faces a different direction. These are small examples of a general problem for situated language agents: communicative success is partner-dependent.

This partner-dependence matters because recent language-agent systems increasingly use interaction as both an evaluation setting and a learning signal. A simulated partner can answer questions, provide feedback, play the role of a user, or supply a reward through task success. Such interaction is attractive for studying embodied and social language use because it moves beyond static question-answer benchmarks and toward multi-turn communication with consequences. But it also creates an evaluation trap. If the partner used during training, filtering, or test-time selection has stable quirks, high interaction reward may mean that the agent has adapted to those quirks rather than learned a robust pragmatic strategy. In other words, an agent can appear socially competent with its training partner while failing with a new partner who interprets ambiguity differently.

This concern echoes a long tradition in pragmatics and coordination. Gricean accounts emphasize informative and cooperative utterance choice~\cite{grice1975logic}; collaborative reference work shows that referring expressions are negotiated with particular partners~\cite{clark1986referring,brennan1996conceptual}; and computational reference-game models formalize how speakers should choose utterances relative to listener alternatives~\cite{frank2012predicting,goodman2016pragmatic}. At the same time, multi-agent learning has shown that self-play can produce conventions that are successful internally but brittle under partner shift, motivating cross-play and zero-shot coordination evaluations~\cite{lewis1969convention,hu2020otherplay,lucas2022anyplay}. Our goal is to connect these views for contemporary language agents: when communicative success is used as a selection signal, evaluation should ask whether success transfers to new listeners.

We study this question in PRAG-CrossPlay, a small, controlled situated reference-game benchmark. A speaker sees a text-rendered grid scene, a hidden target object ID, and the target's visible attributes. The speaker sends a natural-language message, and a listener who does not know the target ID must choose one object from its visible scene. This setting deliberately abstracts away perception in order to isolate pragmatic reference. It also avoids a common source of ambiguity in social-agent evaluation: success is not assigned by a model judge, but by checking whether the listener's chosen object ID matches the ground-truth target.

The key experimental distinction is between \emph{same-play} and \emph{cross-play}. Same-play measures whether a selected message works with the simulated listener used for selection. Cross-play measures whether that same message works with held-out listener models and prompts. We compare a mirror self-play selector, which chooses the shortest candidate that succeeds with one fixed simulated listener, to a population-play selector, which chooses candidates that succeed across a small heterogeneous population of local listeners. We do not claim that this is a full-scale training procedure. Rather, it is a minimal and inspectable instantiation of the broader interaction-learning pattern: a speaker proposes messages, an interaction partner supplies a communicative-success signal, and the system keeps the message that appears to work.

The paper makes three contributions. First, we introduce PRAG-CrossPlay, a released benchmark and protocol with procedural scenes, prompt templates, cached API traces, and verification scripts. Second, we propose a cross-play evaluation protocol for situated natural-language reference that separates same-partner success from held-out communicative success. Third, we provide an empirical diagnosis of partner overfitting in API-based reference games. Mirror self-play can reach perfect same-play success while failing under held-out listeners; population-play closes the observed gap when listener-invariant fallbacks are present; and no-coordinate ablations show that cross-play also exposes weaknesses in the selector itself.

The main empirical picture is intentionally narrow but, we argue, useful. Robust messages are usually available in the candidate pool: oracle selection often reaches perfect held-out success. The failures arise because mirror self-play selects under-informative or frame-sensitive messages that satisfy its fixed partner. This makes PRAG-CrossPlay a compact testbed for a broader claim: communicative success should not be reported only with the same partner that produced the feedback. For situated language agents, cross-partner generalization is a core part of pragmatic evaluation.

\section{Related Work}

\paragraph{Pragmatics and referring expression generation.}
Pragmatic theories of communication emphasize that speakers choose utterances under assumptions about relevance, informativeness, and listener knowledge~\cite{grice1975logic}. Human reference is also interactive: speakers and listeners establish conceptual pacts, adapt lexical choices, and repair ambiguous descriptions over repeated interaction~\cite{clark1986referring,brennan1996conceptual}. In natural-language generation, referring expression generation operationalizes some of these ideas by selecting attributes that distinguish a target from distractors~\cite{dale1995computational}. Rational Speech Act and related probabilistic accounts provide a computational view in which pragmatic speakers choose utterances by reasoning about listener inference~\cite{frank2012predicting,goodman2016pragmatic}. Neural and grounded variants apply these ideas to images, colors, and reference-like settings~\cite{andreas2016reasoning,monroe2017colors,cohngordon2018pragmatically,hu2021scalable}. PRAG-CrossPlay follows this tradition, but asks a different evaluation question: does a message selected by one listener model transfer to other listeners?

\paragraph{Emergent communication, self-play, and cross-play.}
Multi-agent learning has long used self-play and communicative reward to induce conventions or protocols~\cite{foerster2016learning,lazaridou2017multiagent}. The benefit of self-play is that it supplies dense interaction data without human annotation. The risk is that agents may converge to arbitrary conventions that do not transfer to independently trained partners. Zero-shot coordination work addresses this risk by evaluating agents with unseen partners or by designing objectives that reduce convention overfitting~\cite{carroll2019utility,hu2020otherplay,lucas2022anyplay}. Our work imports this diagnostic into natural-language pragmatics. Instead of measuring reward in a cooperative game, we measure whether a natural-language referring expression lets a held-out listener recover a ground-truth referent.

\paragraph{Interactive language-agent evaluation.}
Recent social-agent benchmarks evaluate language models in richer role-play and multi-agent settings. SOTOPIA evaluates agents across open-ended social scenarios with goals such as coordination, exchange, and competition~\cite{zhou2023sotopia}. SOTOPIA-$\pi$ extends this line to interactive learning of socially intelligent language agents and raises the concern that LLM-based evaluators may overestimate agents trained through interaction~\cite{wang2024sotopiapi}. General LLM-as-judge work similarly shows that model-based evaluations can be useful but require care~\cite{zheng2023judging}. PRAG-CrossPlay is narrower by design. It does not aim to measure broad social intelligence; instead, it gives a small objective test of one communicative skill and makes partner overfitting directly measurable.
